{\scriptsize
        \begin{tabularx}{\textwidth}{|p{0.25\columnwidth}|X|}
            \hline 
            \textbf{Notions et contenus} & \textbf{Capacités exigibles} \\
            \hline 
            \multicolumn{2}{|l|}{\textbf{6.2.1. Dispersion et absorption}} \\
            \hline
            Propagation unidimensionnelle d'une onde harmonique dans un milieu linéaire.  & 
            Identifier une analogie avec un phénomène de diffusion.
            Identifier le caractère linéaire d'une équation aux dérivées partielles.
            Établir la relation de dispersion caractéristique d'un phénomène de propagation en utilisant des ondes de la forme $\exp \pm j(k x-\omega t)$.
            Distinguer différents types de comportements selon la valeur de la pulsation. \\
            \hline 
            Dispersion, absorption.   & 
            Associer les parties réelle et imaginaire de $k$ aux phénomènes de dispersion et d'absorption. \\
        \hline 
            Propagation d'un paquet d'ondes dans un milieu non absorbant et faiblement dispersif : vitesse de phase et vitesse de groupe. & 
            Énoncer et exploiter la relation entre les ordres de grandeur de la durée temporelle d'un paquet d'onde et la largeur fréquentielle de son spectre. Déterminer la vitesse de groupe d'un paquet d'ondes à partir de la relation de dispersion.
             Associer la vitesse de groupe à la propagation de l'enveloppe du paquet d'ondes. \\
        \hline
        \end{tabularx}
%         \begin{tabularx}{0.55\textwidth}{|p{0.25\columnwidth}|X|}
%             \hline 
%             \textbf{Notions et contenus} & \textbf{Capacités exigibles} \\
%             \hline 
%             \multicolumn{2}{|l|}{\textbf{6.3. Interfaces entre deux milieux}} \\
%             \hline
%             Réflexion et transmission d'une onde acoustique plane progressive sous incidence normale sur une interface plane infinie entre deux fluides: coefficients de réflexion et de transmission en amplitude des vitesses, des surpressions et des puissances acoustiques surfaciques moyennes.  & 
%             Expliciter des conditions aux limites à une interface.
% Établir les expressions des coefficients de transmission et de réflexion.
% Associer l'adaptation des impédances au transfert maximum de puissance. \\
%             \hline 
%             Réflexion et transmission d'une onde électromagnétique plane progressive harmonique polarisée rectilignement à l'interface entre deux milieux d'indices complexes $\underline{n}_1$ et $\underline{n}_2$ dans le cas d'une incidence normale : coefficients de réflexion et de transmission du champ électrique.   & 
%             Exploiter la continuité admise du champ électromagnétique dans cette configuration pour obtenir l'expression des coefficients de réflexion et de transmission en fonction des indices complexes.
% Utiliser les expressions des coefficients de réflexion et de transmission du champ électrique dans des situations variées. Établir et interpréter les expressions des coefficients de réflexion et de transmission en puissance dans le cas d'une interface entre deux milieux diélectriques linéaires, homogènes, isotropes et transparents. \\
%         \hline
%         \end{tabularx}
\vspace{0.5cm}
}
